% paper/sections/04f_uccd6.tex
%
% §4.6  UCCD6：6 次節点中心 CCD ＋ 次数保存ハイパー粘性
%
% ═══════════════════════════════════════════════════════════════════════════════
\subsection{\texorpdfstring{{\UCCD[6]}：次数保存ハイパー粘性付き節点中心 CCD}{UCCD6：次数保存ハイパー粘性付き節点中心 CCD}}
\label{sec:uccd6_def}
% ═══════════════════════════════════════════════════════════════════════════════

{\UCCD[6]} は，Chu--Fan の CCD 演算子 $D_1^{\CCD}$ の
\textbf{6 次分散精度を保ったまま}
高次ハイパー粘性 $\sigma|a|h^7(-D_2^{\CCD})^4$ を
演算子内部に埋め込むことで，
節点中心のままモーメント移流の $L^2$ 安定性を得る設計である．
DCCD（§\ref{sec:dccd_derivation}，後フィルタ）と
FCCD（§\ref{sec:fccd_def}，面中心）に対する
\textbf{第 3 の経路}として，本節では定式化・安定性解析・境界閉包を述べる．

% -------------------------------------------------------------------------------
\subsubsection{演算子定義と符号}
\label{sec:uccd6_hyperviscosity}
% -------------------------------------------------------------------------------

線形移流 $\partial_t u + a\,\partial_x u = 0$ に対する半離散演算子：
\begin{equation}
  \partial_t U_j
  \;=\; L_h U_j
  \;=\; -a\,(D_1^{\CCD}U)_j \;-\; \sigma|a|h^{7}\,\bigl((-D_2^{\CCD})^{4}U\bigr)_j,
  \qquad \sigma > 0.
  \label{eq:uccd6_semidiscrete}
\end{equation}
ここで $D_1^{\CCD}, D_2^{\CCD}$ は Chu--Fan CCD の 1 次・2 次導関数演算子
（§\ref{sec:ccd_def}）．第二項が\textbf{ハイパー粘性}であり，
$(-D_2^{\CCD})$ が正定値（§\ref{sec:ccd_omega_weights}）であることから
その 4 乗も正定値で，式~\eqref{eq:uccd6_semidiscrete} は\textbf{散逸演算子}を加えた形となる．

指数 7 の必然性を確認する．周期格子での Fourier シンボル：
\begin{equation}
  \lambda(\theta)
  \;=\; -i\,\frac{a\,\omega_1(\theta)}{h}
  \;-\; \sigma|a|\,\frac{\omega_2(\theta)^{8}}{h},
  \qquad \theta = kh,
  \label{eq:uccd6_symbol}
\end{equation}
を用いると，$\omega_2^2 = \theta^2 + \Ord{\theta^8}$
（式~\eqref{eq:ccd_omega_12} 導出系）より
ハイパー粘性は物理波数で
$-\sigma|a|h^{7}\partial_x^{8}u$ として作用し，
主項 $D_1^{\CCD}$ の $\Ord{h^6}$ 切断誤差より\textbf{漸近的に劣位}．
指数 5（$h^{5}(-D_2^{\CCD})^3$）では散逸項が $\omega_2^6/h = \theta^6/h + \Ord{\theta^{12}/h}$
として作用し，主項 $D_1^{\CCD}$ の $\Ord{h^6}$ 切断誤差（$\theta - \theta^7/9450 + \Ord{\theta^9}$）と
低波数で同オーダー競合し精度を削る；
指数 9 では $\theta=\pi$ 近傍モードを減衰しきれない；
指数 7・ハイパー粘性 4 乗は\textbf{次数保存の最小構成}である．

% -------------------------------------------------------------------------------
\subsubsection{半離散 $L^2$ 安定性（エネルギー法）}
\label{sec:uccd6_l2_stability}
% -------------------------------------------------------------------------------

周期格子の $\ell^2$ 内積 $\langle\cdot,\cdot\rangle_h$ の下で
$D_1^{\CCD}$ は歪 Hermite，$D_2^{\CCD}$ は Hermite かつ負定値．従って：
\begin{equation}
  \frac{1}{2}\frac{\mathrm{d}}{\mathrm{d}t}\|U\|_h^{2}
  \;=\; \langle U, \dot U\rangle_h
  \;=\; -\sigma|a|h^{7}\,\bigl\|(-D_2^{\CCD})^{2}U\bigr\|_h^{2}
  \;\le\; 0.
  \label{eq:uccd6_energy_id}
\end{equation}
移流項 $-a\langle U, D_1^{\CCD}U\rangle_h$ は歪 Hermite 性より消滅；
ハイパー粘性項が\textbf{厳密な散逸同定}を与える．
従って UCCD6 は全ての $\theta\in(0,\pi]$ で
モードごとに散逸する（$\omega_2(\theta)\ne 0$）．

% -------------------------------------------------------------------------------
\subsubsection{Crank--Nicolson による無条件安定}
\label{sec:uccd6_cn_stability}
% -------------------------------------------------------------------------------

時間刻み $\tau$ の Crank--Nicolson：
\begin{equation}
  (I + \tfrac{\tau}{2}L_h)\,U^{n+1}
  \;=\; (I - \tfrac{\tau}{2}L_h)\,U^{n}.
  \label{eq:uccd6_cn}
\end{equation}
Fourier 基底で増幅因子：
\begin{equation}
  G(\theta)
  \;=\; \frac{1 + \tau\lambda(\theta)/2}{1 - \tau\lambda(\theta)/2},
  \qquad
  |G(\theta)|^{2}
  \;=\; \frac{(1 - \tau|\Re\lambda|/2)^2 + (\tau\,\Im\lambda/2)^2}{(1 + \tau|\Re\lambda|/2)^2 + (\tau\,\Im\lambda/2)^2}.
  \label{eq:uccd6_amp_factor}
\end{equation}
ここで $\Re\lambda(\theta)\le 0$（式~\eqref{eq:uccd6_symbol}）より
$1 + \tau\Re\lambda/2 = 1 - \tau|\Re\lambda|/2$ および
$1 - \tau\Re\lambda/2 = 1 + \tau|\Re\lambda|/2$ を用いた．
分子 $\le$ 分母より $|G(\theta)|\le 1$ が
$\tau, h$ によらず成立；すなわち UCCD6+CN は\textbf{無条件安定}．

\begin{remark}[陽解法の CFL 拘束]
\label{rem:uccd6_explicit_cfl}
TVD-RK3 などの陽解法を用いると $\omega_2^8$ 依存が
$\theta=\pi$ 近傍で $\tau_{\max}\le \Ord{h/\omega_{2,\max}^8}$ の厳しい拘束を課す．
中程度の $\sigma$ でも実質使えず，CN との組合せが必然．
\end{remark}

% -------------------------------------------------------------------------------
\subsubsection{境界閉包と GKS 文脈}
\label{sec:uccd6_boundary}
% -------------------------------------------------------------------------------

有界区間 $[0,L]$ における Chu--Fan 閉包は
$j=1, N$ で 4 次精度の Hermite 閉包（§\ref{sec:ccd_bc}）を継承する．
$(-D_2^{\CCD})$ の 4 回適用では境界付近で $\Ord{h^4}$ となるが，
$h^7$ 係数がこれを主項 $\Ord{h^6}$ より抑制するため大局的精度は保たれる．
厳密な Gustafsson--Kreiss--Sundstr\"om 安定性は今後の課題とする．

% -------------------------------------------------------------------------------
\subsubsection{3 スキームの役割分担}
\label{sec:uccd6_scheme_comparison}
% -------------------------------------------------------------------------------

\begin{center}
\begin{tabular}{l l l l}
\hline
\textbf{スキーム} & \textbf{配置} & \textbf{散逸チャンネル} & \textbf{保存精度} \\ \hline
Chu--Fan CCD（§\ref{sec:ccd_def}） & 節点 & なし & $\Ord{h^6}$ \\
DCCD（§\ref{sec:dccd_derivation}） & 節点 & 外殻 post-filter & $\Ord{h^6}$ \\
FCCD（§\ref{sec:fccd_def}） & \textbf{面} & 風上構造 & $\Ord{h^4}$ \\
\textbf{UCCD6（本節）} & 節点 & \textbf{演算子内部} & $\Ord{h^6}$ \\
\hline
\end{tabular}
\end{center}

\noindent
本稿では UCCD6 が bulk 移流の既定として DCCD を置換する．
$\omega_2^8$ の選択的 Nyquist 散逸は DCCD の一様減衰 $H(\pi)=1-4\varepsilon_d$ より
dispersion preservation に優り，bulk 高周波の保持性で上位に位置する．
DCCD は UCCD6 の動機付けを成す教育的前段（§\ref{sec:dccd_derivation}）として
および CLS 再初期化圧縮項（§\ref{sec:reinit}，非線形フラックス散逸子）として保持される．

\noindent
UCCD6 は FCCD と\textbf{直交補完}の関係にある：
FCCD の $\Ord{H^4}$ は界面消費点（(R1)-(R4)；§\ref{sec:fccd_def}）のみに展開され，
バルク演算子は節点 CCD $\Ord{h^6}$ を維持する（設計選択であり精度の妥協ではない）．
UCCD6 は節点配置のまま演算子内部散逸でモーメント移流の Gibbs を抑制し，
Level-1 検証・Level-2 製品・Level-3 DNS すべての bulk advection 既定として
§\ref{sec:time_int} の時間積分スタックに組み込まれる．

% ═══════════════════════════════════════════════════════════════════════════════
